\documentclass{article}
\usepackage[utf8]{inputenc}
\usepackage[T1]{fontenc}
\usepackage{geometry}
\usepackage{graphicx}
\usepackage{float}
\usepackage{hyperref}
\usepackage[english,indonesian]{babel}

\title{Optimasi Pemrosesan Template LaTeX Menggunakan Pendekatan Berbasis AI}
\author{Academic Support Agent}
\date{\today}

\begin{document}
\maketitle

\begin{abstract}

\end{abstract}

\chapter{PENDAHULUAN}


\chapter{TINJAUAN PUSTAKA}


% Original content (commented out):
% \documentclass[12pt,a4paper]{article} % 设置文档类型为article，基本字体大小为12pt，纸张大小为A4
\usepackage[utf8]{inputenc} % 使用UTF-8编码，以支持多语言文本输入
\usepackage[T1]{fontenc} % 设置字体编码为T1
\usepackage{ctex} % 引入ctex包，以支持中文排版
\usepackage[left=25mm,right=25mm,top=25mm,bottom=25mm,headheight=15mm,footskip=15mm]{geometry} % 设置页边距和页眉、页脚高度
\usepackage{fancyhdr} % 使用fancyhdr包来自定义页眉和页脚
\usepackage{titlesec} % 使用titlesec包来自定义章节标题的格式
\usepackage{setspace} % 引入setspace包，用于设置行间距
\usepackage[titletoc,title]{appendix} % 使用appendix包来自定义附录部分
\usepackage{lipsum} % 引入lipsum包，用于生成示例文本
\usepackage{times} % 使用Times New Roman字体


\pagestyle{fancy} % 使用fancy样式
\fancyhf{} % 清空当前的页眉和页脚设置
\fancyhead[R]{\songti\small\textit{湖北大学本科毕业论文（设计）}} % 居右
\fancyfoot[C]{\thepage} % 设置页脚中间为页码
\renewcommand{\headrulewidth}{0pt} % 移除页眉下方的横线

% 页码样式设置，这里的代码需根据上下文具体问题具体分析，可能需调整
\renewcommand{\thepage}{\ifnum\value{page}<1\relax\Roman{page}\else\arabic{page}\fi}
\renewcommand{\contentsname}{\heiti\zihao{3}\centering 目\hspace{1em}录}


% 设置段落首行缩进和段落间距
\setlength{\parindent}{2em} % 设置首行缩进为2个字符
\setlength{\parskip}{0em} % 设置段落间距为0
\linespread{1} % 设置行间距为单倍

\begin{document}

\input{chapter/cover}


% 封面、目录等前置部分使用罗马数字页码
\pagenumbering{roman} % 将页码设置为罗马数字


%目录
\newpage
\tableofcontents

% 中文摘要页
\input{chapter/chapter1-ChineseAbstract}

% 英文摘要页
\input{chapter/chapter2-EnglishAbstract}

% 中文部分
\input{chapter/chapter3-text}

% 引用部分
\input{chapter/chapter4-reference}

% 附录
\input{chapter/chapter5-appendix}

% 致谢
\input{chapter/chapter6-thanks}

\end{document}

\end{document}